\section{Background}
\label{sec:background}

\paragraph{Filament motion.}
A filament is a field-aligned tube of plasma, with a cross-field radius $a$ between a few ion gyroradii and a few centimetres, that is denser and hotter than the surrounding SOL. In the outboard SOL the magnetic curvature and gradient drifts separate ions and electrons across the filament, and the resulting electric field drives an $E \times B$ drift that moves the whole filament radially outward \citep{krasheninnikov2001scrape, dippolito2011convective}. The drift speed depends on the potential that the charge separation can build, and therefore on the paths through which the polarization charge is dissipated, which leads to two standard limits. If the filament is connected along the field line to the divertor target, the charge drains through the sheath and $v \propto c_s (L_\parallel/R)(\rho_s/a)^2$, where $c_s$ is the sound speed, $\rho_s$ the ion sound gyroradius, $L_\parallel$ the connection length, and $R$ the major radius \citep{krasheninnikov2001scrape}, whereas if the ion polarization current closes the circuit inside the filament, $v \propto c_s (a/R)^{1/2}$ \citep{garcia2006fluctuations}.

\paragraph{The two-region model.}
\citet{myra2006collisionality} model the midplane and the divertor leg as two coupled regions and join these limits. With
\begin{equation}
    a^* = \rho_s^{4/5} L_\parallel^{2/5} R^{-1/5}, \qquad
    v^* = c_s \left(\frac{2 a^*}{R}\right)^{1/2}, \qquad
    \Theta = \left(\frac{a}{a^*}\right)^{5/2}, \qquad
    \Lambda = \frac{\nu_{ei} L_\parallel}{\Omega_e \rho_s},
    \label{eq:tworegion}
\end{equation}
and $\hat a = a/a^*$, $\hat v = v/v^*$, the model gives
\begin{equation}
    \hat v =
    \begin{cases}
        \hat a^{1/2} & \text{connected ideal interchange } (\Lambda < 1,\ \Theta < 1), \\
        \hat a^{-2} & \text{sheath connected } (\Lambda < 1,\ \Theta > 1), \\
        \Lambda\, \hat a^{-2} & \text{resistive X-point } (1 < \Lambda < \Theta), \\
        \hat a^{1/2} & \text{resistive ballooning } (\Lambda > 1,\ \Lambda > \Theta).
    \end{cases}
    \label{eq:regimes}
\end{equation}
Here $\nu_{ei}$ is the electron--ion collision frequency and $\Omega_e$ the electron gyrofrequency \citep{dippolito2011convective}, and at low collisionality the velocity rises as $\hat a^{1/2}$ for small filaments, peaks near $\hat a = 1$, and falls as $\hat a^{-2}$ for large ones. The model is a scaling argument whose prefactors depend on the filament profile and on the parallel structure, and simulations in realistic geometry find velocities below the model values by factors of order unity \citep{easy2014three, walkden2015characterization, paruta2018blob}. \citet{theiler2009blob} measured the velocity--size relation on the TORPEX basic plasma device and recovered the predicted transitions between regimes.

\paragraph{Camera measurements.}
Visible-light imaging records line-integrated $D_\alpha$ emission, which depends on the local density, temperature, and neutral density. Filaments appear as bright, field-aligned streaks against the quiescent background, and their emission tracks the density perturbation closely enough that velocities and widths from imaging agree with probe measurements within tens of percent \citep{zweben2007edge, benayed2009interelm}. In a tangential view the brightest image feature is the plasma limb, the band where the line of sight is tangent to the flux surfaces near the LCFS, and filaments cross the limb as they move outward. Because the camera integrates over each exposure, a filament that moves a distance $v\,t_{\rm exp}$ during the exposure is smeared by that distance along its direction of motion, an effect treated explicitly in Section~\ref{sec:analysis}.
